\documentclass[12pt]{article}

\usepackage{amsmath,amssymb,amsthm,hyperref,geometry}
\geometry{margin=1in}

\title{\bf Ehlers Frame Theory and the Newtonian Limit of Relativistic Gravitation}
\author{ }
\date{ }

\begin{document}
\maketitle

\begin{abstract}
Ehlers frame theory provides a unified, covariant framework that treats
general relativity and Newton–Cartan gravity as members of a single
one-parameter family of spacetime theories. 
Contrary to common simplifications, the non-relativistic limit of general
relativity cannot be obtained by any naive procedure such as $c \to \infty$.
Instead, Newtonian gravity appears as a geometric limit of a family of 
Lorentzian spacetimes only under strict structural and topological conditions.
This article presents a modern, detailed exposition of Ehlers frame theory,
its mathematical axioms, limiting procedures, and implications for the
structure of classical mechanics in the absence of global time.
We further discuss connections to ADM theory, global hyperbolicity, and the 
geometry of temporal and spatial metric structures.
\end{abstract}

\tableofcontents

\section{Introduction}

The problem of understanding the Newtonian limit of general relativity is
conceptually subtle.
Newtonian gravity is formulated on a manifold equipped with a degenerate pair
of temporal and spatial metric structures, $(t_a, h^{ab})$, whereas general
relativity uses a non-degenerate Lorentzian metric $g_{ab}$.
The two theories live on geometrically incompatible structures, and therefore
one cannot obtain Newton–Cartan theory through a simple limit $c\to\infty$.

Ehlers frame theory, introduced by J.~Ehlers in a series of investigations on
the structure of gravitation, provides an elegant solution: construct a
one-parameter family of geometric structures $(t_{ab}(\lambda),
h^{ab}(\lambda),\nabla(\lambda))$, $\lambda \ge 0$, such that for $\lambda>0$
these fields describe a relativistic spacetime with Lorentzian metric, while
as $\lambda \to 0$ one obtains a Newton–Cartan spacetime.
This approach allows for precise, covariant statements about when a Newtonian
limit exists, what geometric structures survive the limit, and how matter and
gravitational fields must scale in the non-relativistic regime.

In addition to clarifying the classical limit, frame theory also offers
insight into why classical mechanics fails when no global time function exists.
In relativistic theory, foliations into spatial hypersurfaces are not generally
available, so notions such as rigid body, global simultaneity, or canonical
phase space structures may fail to exist. 
Newtonian physics is recovered only when the temporal metric structure becomes
integrable in the $\lambda\to 0$ regime.

This paper aims to present a long, journal-style exposition of the frame theory
framework, suitable for readers familiar with differential geometry,
relativistic field theory, and Newton–Cartan geometry.

\section{The Structural Axioms of Frame Theory}

A \emph{frame theory} consists of the following data on a smooth, connected
four-dimensional manifold $M$.

\begin{enumerate}
\item A parameter $\lambda \ge 0$, called the \emph{causality parameter}.
\item A symmetric covariant tensor $t_{ab}(\lambda)$ (temporal metric).
\item A symmetric contravariant tensor $h^{ab}(\lambda)$ (spatial metric).
\item A torsion-free affine connection $\nabla(\lambda)$.
\item Matter fields described by an EFT stress-energy tensor $T^{ab}(\lambda)$.
\end{enumerate}

These fields satisfy the fundamental \emph{orthogonality relation}
\begin{equation}
  t_{ab}(\lambda) \, h^{bc}(\lambda) = - \lambda \, \delta_a^c.
  \label{orth}
\end{equation}
The parameter $\lambda$ controls the causal structure. 
For $\lambda>0$, $t_{ab}$ and $h^{ab}$ combine into a non-degenerate metric.
For $\lambda=0$, the right-hand side vanishes and $t_{ab}h^{bc}=0$, which
forces both tensors to be degenerate and orthogonal, as in Newton–Cartan
geometry.

The connection is compatible with these structures:
\begin{equation}
  \nabla_a t_{bc} = 0, \qquad 
  \nabla_a h^{bc} = 0.
  \label{compat}
\end{equation}
These conditions generalize the Levi--Civita compatibility conditions.
When $\lambda>0$, the orthogonality relation ensures that the combined metric
\begin{equation}
  g_{ab}(\lambda) 
    := t_{ab}(\lambda) + \frac{1}{\lambda} h_{ab}(\lambda),
  \label{metric}
\end{equation}
where $h_{ab}$ is the inverse of $h^{ab}$ on the spatial subspace,
is a Lorentzian metric, and \eqref{compat} ensures that
$\nabla(\lambda)$ is its Levi--Civita connection.

Thus for every positive $\lambda$ the theory reduces to ordinary general
relativity, provided one also imposes the Einstein field equations
\begin{equation}
  G_{ab}(g) = 8\pi T_{ab}.
\end{equation}

\section{The Limit $\lambda\to 0$ and Newton--Cartan Geometry}

The key insight of Ehlers is that Newtonian gravity emerges only when one
considers a \emph{family} of relativistic spacetimes whose temporal and spatial
metrics scale such that \eqref{orth} has a smooth limit.

As $\lambda\to 0$, the tensors become degenerate:
\[
\mathrm{rank}( t_{ab}(0) ) = 1, 
\qquad 
\mathrm{rank}( h^{ab}(0) ) = 3.
\]
Their null spaces define spatial and temporal subspaces.

Let $t_a := t_{ab}(0)$ denote the absolute time one-form, uniquely defined up
to scale if integrability conditions hold.
The spatial metric $h^{ab}(0)$ is the inverse of the three-dimensional spatial
metric on the fibers of $\ker t_a$.

The limiting connection $\nabla(0)$ is not Levi--Civita, since no non-degenerate
metric exists.
Instead, it is a Newton--Cartan connection, satisfying
\[
\nabla_a t_b = 0, \qquad \nabla_a h^{bc} = 0,
\]
and having curvature components associated with the Newtonian gravitational
field.
When matter is included, the temporal-temporal curvature reproduces Poisson
gravity,
\begin{equation}
  h^{ij} R_{i0j0} = \nabla^2 \Phi,
\end{equation}
with $\Phi$ the Newtonian potential.
Thus Newtonian gravity is not simply a special case of relativity but the
geometric limit of a specific class of relativistic spacetimes.

\section{Existence and Obstructions of the Newtonian Limit}

Not every relativistic spacetime can appear as part of a frame theory
approaching a Newton--Cartan structure.

\subsection{Topological and global requirements}

A necessary condition is the existence of a global time function 
$\tau: M \to \mathbb{R}$ whose level sets are spacelike.
Spacetimes without a global foliation by spacelike hypersurfaces—such as Gödel
spacetime or general rotating cosmologies—do not admit a Newtonian limit in the
Ehlers sense.
This demonstrates that Newtonian theory implicitly assumes strong global
structure: absolute simultaneity requires integrability of the temporal 
one-form $t_a$.

\subsection{Curvature scaling}

Curvature must scale in a controlled way:
\[
|R^a{}_{bcd}(\lambda)| < C
\]
with $\lambda$-dependent contractions respecting the scaling of $t_{ab}$ and
$h^{ab}$.
For example, one requires that the ``spatial'' curvature components remain
finite, while ``mixed'' components scale appropriately to produce the
Newtonian tidal field.

\section{Relation to ADM Splitting and Temporal Functions}

ADM theory provides a foliation-based decomposition of Lorentzian spacetimes
into space-plus-time.
However, ADM requires a global time function whose level sets are spacelike.
Ehlers frame theory clarifies which aspects of ADM survive in the 
$\lambda\to 0$ limit.

If $(M,g_{ab})$ admits a global time function $t$ with timelike gradient, one
can express the metric as
\[
g_{ab} = -N^2 t_a t_b + h_{ab},
\]
where $h_{ab}$ is the spatial metric induced on slices $t=\mathrm{const}$.
In the frame theory scaling, the lapse $N$ corresponds to the clock metric
encoded in $t_{ab}(\lambda)$.
As $\lambda\to 0$, the spatial metric $h_{ab}$ becomes the inverse of
$h^{ab}(0)$, and the extrinsic curvature components converge to the spatial
connection of Newton--Cartan theory.

\section{Implications for Classical Mechanics and Rigid Bodies}

Classical mechanics requires:
\begin{itemize}
\item a global time,
\item spatial slices that define positions of bodies at each instant,
\item a tangent bundle splitting into temporal and spatial directions.
\end{itemize}

Frame theory shows that these structures exist only when
$t_{ab}(\lambda)$ becomes integrable as $\lambda\to0$.
If a relativistic spacetime does not admit a global time function,
Newtonian mechanics cannot be recovered.

This explains the non-existence of rigid bodies and holonomic constraints in
relativistic physics: spatial simultaneity and rigid spatial distances require
a foliation and spatial metric independent of temporal structure, both of which
are lost in general relativity unless special symmetries exist.

\section{Matter Coupling in Frame Theory}

Matter is described by a tensor $T^{ab}(\lambda)$ satisfying
\[
\nabla_a T^{ab} = 0.
\]
The Newtonian limit imposes scaling conditions, such as:
\[
T^{00}(\lambda) \sim O(1), \quad 
T^{0i}(\lambda) \sim O(\sqrt{\lambda}), \quad
T^{ij}(\lambda) \sim O(\lambda).
\]
These scalings ensure the emergence of Newtonian mass density and momentum
current, and eliminate relativistic stresses when $\lambda=0$.

\section{Conclusions}

Ehlers frame theory provides a mathematically rigorous bridge between general
relativity and Newtonian gravity.
It clarifies that Newtonian theory is not simply a low-velocity approximation
but a singular geometric limit requiring strong global assumptions, including
the existence of a global time function and appropriate curvature scaling.
The theory offers insight into the deep structural incompatibilities between
classical and relativistic mechanics, and gives a natural home for
post-Newtonian approximations, ADM initial-value formulations, and the geometry
of absolute time.

\end{document}

